% Generated draft: requires mathematical and source review.
\documentclass[11pt]{article}
\usepackage[margin=25mm]{geometry}
\usepackage{fontspec}
\setmainfont{Latin Modern Roman}
\usepackage{amsmath,amssymb,mathtools,unicode-math}
\setmathfont{Latin Modern Math}
\usepackage{enumitem,needspace,xurl,hyperref}
\setlength{\parindent}{0pt}
\setlength{\parskip}{5pt}
\setlength{\emergencystretch}{3em}
\newcommand{\N}{\mathbb{N}}
\newcommand{\Z}{\mathbb{Z}}
\newcommand{\Q}{\mathbb{Q}}
\newcommand{\R}{\mathbb{R}}
\newcommand{\C}{\mathbb{C}}
\newcommand{\F}{\mathbb{F}}
\newcommand{\A}{\mathbb{A}}
\newcommand{\PP}{\mathbb{P}}
\newcommand{\E}{\mathbb{E}}
\newcommand{\eps}{\varepsilon}
\newcommand{\dd}{\,\mathrm{d}}
\begin{document}
\section*{UnsolvedMath ID: 3700009}
\textbf{Source problem code:} catalogue entry
\subsection{Definitions and mathematical statement}
The source concerns a rational map \(f:\widehat{\mathbb C}\to\widehat{\mathbb C}\) and an analytic invariant curve \(C\). In the usual setting of rational dynamics one takes \(\deg f\ge2\); the source excerpt does not expressly state this bound or the precise class of curves. Let \(J(f)\) denote the Julia set, the complement of the open set where the iterates of \(f\) form a normal family. A fixed point \(z\) is repelling when the absolute value of its multiplier in a local coordinate is greater than \(1\). A critical point is a point where the local mapping degree exceeds \(1\).

Retaining the source's wording, label its additional conditions as follows:
\[
\begin{array}{ll}
H_1:& f|_C:C\to C\text{ is not a homeomorphism};\\
H_2:& C\text{ contains a repelling fixed point of }f;\\
H_3:& C\subseteq J(f);\\
H_4:& C\text{ contains no critical point of }f;\\
H_5:& C\text{ contains no ``rational fixed point'' of }f.
\end{array}
\]
The source reports a theorem with conclusion
\[
 f\text{ is a Latt\`es function}\quad\text{or}\quad C\text{ is algebraic},
\]
and asks which of these hypotheses may be omitted while that conclusion remains true for every admissible \((f,C)\). More explicitly, for a proposed set of retained conditions \(I\subseteq\{1,2,3,4,5\}\), the intended test is whether every rational map and analytic invariant curve in the theorem's underlying class satisfying \(H_i\) for \(i\in I\) still satisfies the dichotomy. The source does not specify whether removal of the underlying analytic-curve or invariance assumptions is also being contemplated.

Several terms prevent the excerpt alone from being a complete theorem statement. It does not say whether ``analytic curve'' means a real-analytic Jordan curve, an arc, or another class, nor whether invariance means \(f(C)=C\) or merely \(f(C)\subseteq C\). It gives no definition of ``rational fixed point''; this must not be silently replaced by ``rationally indifferent fixed point'' or by a coordinate-arithmetic condition. The meaning of ``algebraic'' must be fixed for the intended curve class. The customary Latt\`es terminology concerns a rational map induced by an affine map on a complex torus through a finite holomorphic quotient, but the original theorem's precise convention is not supplied.

The research summary says the dichotomy is proved under hypotheses ``(a) or (b)'', whereas neither alternative is included or identified with the displayed list in the statement. Thus the exact established theorem and the allowed removals require the original source. The supplied classification is partial progress, with sharpness open; mentions of M\"uller and Pakovich's semiconjugacy-factorization work are reported as relevant, not as resolving these removals.
\subsection{Short English statement}
For a rational map preserving an analytic curve, which assumptions are needed to force the map to be Latt\`es or the curve to be algebraic? The excerpt lists several assumptions but leaves the curve class, ``rational fixed point'', and two referenced alternative hypotheses unspecified, so the exact theorem needs source review.
\subsection{Sources}
\begin{itemize}
\item Local registry: \url{data/unsolvedmath/problems.json}, permanent UnsolvedMath ID \(3700009\), source code catalogue entry; supplied statement, background, and research summary, updated 2026-08-06. Attribution: Alexandre Eremenko and cited proposers; registry proposed year 2025.
\item Alexandre Eremenko, \emph{My favorite unsolved problems}, ``Analytic invariant curves'', Question 3: \url{https://www.math.purdue.edu/~eremenko/uns1.html}. Supplied extraction method: PDF text; accessed 2026-07-29. The record says explicit solved/counterexample updates in the author-maintained index were applied through that access date.
\item The research summary mentions M\"uller and Pakovich's classification of semiconjugacy triples, but provides no paper titles, dates, theorem statements, or direct citations for the reported dichotomy.
\end{itemize}
\end{document}
